\documentclass[11pt,a4paper]{article}

%%%%%%%%%%%%%%%%%%%%
% STANDARD PACKAGES
%%%%%%%%%%%%%%%%%%%%
\usepackage[utf8]{inputenc}
\usepackage[T1]{fontenc}
\usepackage{lmodern}
\usepackage[margin=1in]{geometry}
\usepackage{amsmath,amssymb,amsfonts,bm}
\usepackage{siunitx}
\usepackage{physics}
\usepackage{graphicx}
\usepackage{booktabs}
\usepackage[numbers,super,sort&compress]{natbib}
\usepackage[colorlinks=true,allcolors=blue]{hyperref}

% Fix for siunitx and physics package conflict
\AtBeginDocument{\RenewCommandCopy\qty\SI}

% Custom Command for Covariant D'Alembertian
\newcommand{\wave}{\nabla^\mu\nabla_\mu}

%%%%%%%%%%%%%%%%%%%%
% TITLE & AUTHOR
%%%%%%%%%%%%%%%%%%%%
\title{\textbf{A Mechanical Route to \texorpdfstring{$G$}{G}: Matching Field Stress--Energy to the Einstein--Hilbert Coupling}}

\author{
    Robin Victor Skavberg\thanks{ORCID: \href{https://orcid.org/0009-0003-9126-6196}{0009-0003-9126-6196}} \\
    \small Independent Researcher \\
    \small Calgary, Alberta, Canada
}

\date{January 26, 2026}


%%%%%%%%%%%%%%%%%%%%
% BEGIN DOCUMENT
%%%%%%%%%%%%%%%%%%%%
\begin{document}

\maketitle

\begin{abstract}
We develop a mechanical route to Newton’s gravitational coupling $G$ by relating the macroscopic field stress--energy to the geometric stiffness of the Einstein--Hilbert (EH) action. In the long--wavelength regime of a homogeneous, barotropic background, the bulk modulus satisfies $K=\rho_m c_s^2$; identifying the signal speed with $c$ yields $K=\rho_m c^2$. We then show that an EH term $\int\!\sqrt{-g}\,R$ arises with a calculable coefficient in the one--loop effective action, so that matching the geometric stiffness $c^4/(8\pi G)$ to the mechanical stiffness gives $G=c^2/(8\pi\rho_m)$. Independently, we verify the observational identity $G=\Lambda c^4/(8\pi u_\nu)$, where $\Lambda$ is the cosmological constant and $u_\nu$ is the corresponding vacuum stress--energy density, providing a parameter--free cross--check.

To connect with microphysics and phenomenology, we outline three realizations: (A) an elastic curvature response consistent with thermodynamic and induced--gravity perspectives; (B) a superfluid dark--matter realization in which phonon excitations couple to baryons and exert a literal force; and (C) a single--field EFT of dark energy with explicit matter couplings that generate an effective force and $G_{\rm eff}(t,k)$. Our baseline adopts (A) (no extra force) and recovers $\Lambda$CDM at linear order; (B) and (C) are presented as optional model ingredients with distinct, testable signatures.
\end{abstract}


In a zero-sum balanced universe, $G$ emerges as the necessary proportionality constant required to maintain the equilibrium between the vacuum's mechanical stiffness and the observed geometric curvature. 
\end{abstract}

\section{Introduction}
Analogue--gravity studies show that hydrodynamic media can reproduce relativistic kinematics via an effective metric,\cite{Unruh1981,Visser1998,Barcelo2011} suggesting a conservative route to relate gravitational dynamics to mechanical properties. In this work, we develop a mechanical route to $G$ by modeling the vacuum as a quantum condensate. Identifying the signal speed with $c$ fixes the field pressure to $K=\rho_m c^2$, which we match to the geometric stiffness multiplying curvature in the Einstein–Hilbert (EH) action.\cite{Einstein1916,Hilbert1915}

\section{Assumptions and Regime of Validity}
We assume: (i) long--wavelength hydrodynamics ($k\xi\ll1$); (ii) homogeneous, barotropic background with mass density $\rho_m$; (iii) phase coherence; (iv) identification of acoustic and electromagnetic null cones in the IR ($c_s=c$).

\section{Field Pressure: Three Realizations and Couplings}
\label{sec:three-realizations}

We use ``field pressure'' as shorthand for the macroscopic stress--energy that sets a response scale for spacetime. The way this response acts on matter depends on microphysics and couplings. We distinguish three realizations:

\paragraph{(A) Elastic curvature response (baseline).}
With minimal coupling and a cosmological constant background, the field pressure fixes the medium's elastic response; matter follows geodesics of the curved metric. The ``push back'' is spacetime curvature---an equation-of-state relation---not a direct pressure force on baryons.\cite{Jacobson1995,Padmanabhan2004,Sakharov1967}

\paragraph{(B) Superfluid dark matter with phonon force (optional).}
If the field is a dark-matter superfluid in galaxies, the phonon mode $\theta$ can couple to baryons, e.g.
$\mathcal{L}_{\rm int}=\alpha\,\theta\,\rho_b$, so baryons feel an additional force
$\mathbf{a}_{\rm phonon}=-\alpha\,\nabla\theta$ on top of geodesic motion. Here pressure/stress gradients transmit a literal force in the superfluid phase, while cosmology remains CDM-like.\cite{BerezhianiKhoury2015,KhouryPLB2015}

\paragraph{(C) Single-field EFT of dark energy with explicit couplings (optional).}
In a general single-field EFT, the background field supplies the pressure/energy density, and matter may couple conformally/disformally to the scalar sector. At the perturbative level this produces an effective force on baryons and a scale- and time-dependent $G_{\rm eff}(t,k)$, subject to stability and observational constraints; setting extra quadratic operators to zero reproduces $\Lambda$CDM at linear order.\cite{Gubitosi2013}

\section{Hydrodynamic Framework}
At the microscopic level, we take a complex scalar condensate described by the Gross--Pitaevskii (GP) equation:\cite{Gross1961,Pitaevskii1961}
\begin{equation}
i\hbar\,\partial_t\Psi=\left[-\frac{\hbar^2}{2m}\nabla^2 + V_{\rm ext} + g|\Psi|^2\right]\Psi.
\end{equation}
Writing $\Psi=\sqrt{n}\,e^{i\theta}$ gives hydrodynamics for number density $n$ and velocity $\vb{v}=(\hbar/m)\nabla\theta$. Linear perturbations propagate on an acoustic metric $ds^2\propto -c_s^2\,dt^2 + d\vb{x}^2$.

\section{Mechanical Stiffness and Geometric Coupling}

\subsection{Bulk Modulus from the Equation of State}
For the GP fluid, the pressure is $P=\tfrac12 g n^2$. We define $\rho_m \equiv mn$ as the mass density of the condensate (where $m$ is the constituent mass and $n$ is the number density). The sound speed is determined by the derivative of pressure with respect to the mass density:\cite{LandauLifshitz}
\begin{equation}
c_s^2 = \frac{dP}{d\rho_m} = \frac{d(\tfrac{1}{2}gn^2)}{d(mn)} = \frac{gn}{m}.
\end{equation}
The bulk modulus (mechanical stiffness) is $\beta \equiv \rho_m \frac{dP}{d\rho_m}$. Setting $c_s=c$ yields:
\begin{equation}
K \equiv \beta = \rho_m c^2.
\end{equation}

\subsection{Matching to Einstein’s Coupling}

The EH action has the prefactor $c^4/(16\pi G)$, corresponding to a geometric stiffness $c^4/(8\pi G)$. 
Equating mechanical and geometric stiffnesses gives:
\begin{equation}
\label{eq:Gfromrho}
\rho_m c^2 = \frac{c^4}{8\pi G}
\;\;\implies\;\;
G = \frac{c^2}{8\pi \rho_m}.
\end{equation}

\subsection{Verification and Significance}
The identity
\begin{equation}
G=\frac{\Lambda c^4}{8\pi u_\nu}
\end{equation}
follows from the standard equivalence between the cosmological constant term
$\Lambda g_{\mu\nu}$ and a Lorentz-invariant vacuum stress--energy
$T^{(\Lambda)}_{\mu\nu}=-u_\nu g_{\mu\nu}$ with $u_\nu=\rho_\Lambda c^2=\Lambda c^4/(8\pi G)$.
Using observational values for $\Lambda$ and $u_\nu$ yields
$G\simeq \SI{6.67e-11}{m^3.kg^{-1}.s^{-2}}$, a parameter-free cross-check.\cite{Jacobson1995}
In our constitutive interpretation, the background vacuum stress--energy (associated with $\Lambda$)
sets the elastic response scale of spacetime; $G$ is the corresponding modulus.
Realizations (B) and (C) introduce explicit matter couplings through which pressure/stress gradients
can exert a literal force on baryons; our baseline analysis adopts (A) with no extra force and hence
reduces to $\Lambda$CDM at linear order.\cite{Gubitosi2013,BerezhianiKhoury2015}

\paragraph{Interpretive consistency.}
Our viewpoint is conceptually consistent with: (i) Jacobson’s reading of the Einstein equations
as an \emph{equation of state}; (ii) Sakharov’s induced-gravity picture in which the
Einstein--Hilbert term encodes a \emph{metric elasticity} generated by vacuum fluctuations; and
(iii) Padmanabhan’s paradigm of \emph{gravity as spacetime elasticity}.
Accordingly, we take the background vacuum stress--energy to set the \emph{macroscopic stiffness}
of the gravitational medium. Newton’s constant $G$ then appears as a \emph{constitutive parameter}
(effective modulus) quantifying spacetime’s resistance to curvature; matter curves spacetime, and
the elastic response is fixed by this modulus.\cite{Jacobson1995,Sakharov1967,Padmanabhan2004}

\subsection{Novel Correlation and Derivation}
It is important to emphasize that while the mathematical components utilized in this derivation—the Gross--Pitaevskii equation, the heat kernel expansion for induced gravity, and the Friedmann cosmological parameters \textit{et al.}—are established in their respective fields, the \textbf{correlation} presented here is unique. We demonstrate that these independent frameworks are not merely parallel but are coupled through a specific mechanical-to-geometric matching condition. By showing that the mechanical bulk modulus of a dark-matter medium is identically equal to the EH prefactor, we provide a new physical derivation for the value of $G$, transforming it from a bare fundamental constant into a calculable property of the vacuum condensate.


While the GPE describes the medium and the Heat Kernel describes the geometry, standard physics keeps them separate. The contribution pertains to the Stiffness Matching Principle. We treat the Einstein-Hilbert action as the 'stress-strain' equation for the vacuum. By equating the mechanical resistance of the dark-matter condensate to the geometric resistance of spacetime, $G$ is no longer a 'given' parameter; it becomes the specific value required for the medium and geometry to remain in equilibrium."

\section{Dimensional Consistency and Physical Interpretation}

To ensure the mathematical validity of the proposed mechanical route, we perform a dimensional analysis of the identity $G = \Lambda c^4 / (8\pi u_\nu)$. As shown in Table~\ref{tab:dimensions}, the units of the derived coupling perfectly match the standard dimensions of Newton's gravitational constant.

\begin{table}[h]
\centering
\caption{Dimensional Analysis of the Mechanical Route to $G$}
\label{tab:dimensions}
\vspace{2mm}
\begin{tabular}{@{}llll@{}}
\toprule
\textbf{Quantity} & \textbf{Symbol} & \textbf{SI Units} & \textbf{Dimensions} \\ \midrule
Gravitational Constant & $G$ & \si{m^3.kg^{-1}.s^{-2}} & $[L]^3 [M]^{-1} [T]^{-2}$ \\
Speed of Light & $c$ & \si{m.s^{-1}} & $[L] [T]^{-1}$ \\
Vacuum Energy Density & $u_\nu$ & \si{kg.m^{-1}.s^{-2}} (Pa) & $[M] [L]^{-1} [T]^{-2}$ \\
Cosmological Constant & $\Lambda$ & \si{m^{-2}} & $[L]^{-2}$ \\
Condensate Density & $\rho_m$ & \si{kg.m^{-3}} & $[M] [L]^{-3}$ \\ \bottomrule
\end{tabular}
\end{table}

\begin{table}[h]
\centering
\caption{Algebraic Breakdown of Dimensional Consistency}
\label{tab:algebraic_dims}
\vspace{2mm}
\begin{tabular}{@{}ll@{}}
\toprule
\textbf{Component} & \textbf{Dimensional Operation} \\ \midrule
Numerator ($[\Lambda] \times [c]^4$) & $[L]^{-2} \times ([L][T]^{-1})^4 = [L]^2 [T]^{-4}$ \\ \addlinespace
Denominator ($[u_\nu]$) & $[M] [L]^{-1} [T]^{-2}$ \\ \addlinespace
\textbf{Final Ratio} ($\frac{\text{Numerator}}{\text{Denominator}}$) & $\frac{[L]^2 [T]^{-4}}{[M] [L]^{-1} [T]^{-2}} = \mathbf{[L]^3 [M]^{-1} [T]^{-2}}$ \\ \midrule
\textbf{Identity Check} & $\equiv [G]$ \\ \bottomrule
\end{tabular}
\end{table}

\subsection*{Physical Significance of the Result}
The dimensional alignment supports the following physical interpretations of the theory:
\begin{itemize}
    \item \textbf{Energy Density as Field Pressure:} The vacuum energy density $u_\nu$ (measured in Pascals) represents the physical "push back" of the dark-matter condensate against gravitational compression.
    \item \textbf{Curvature as Strain:} The cosmological constant $\Lambda$ (measured in inverse area) functions as the intrinsic strain or geometric deformation of the medium over cosmological scales.
    \item \textbf{Gravity as an Emergent Stiffness:} The gravitational constant $G$ is the resulting modulus that scales the relationship between energy density and geometry, functioning as the emergent "stiffness constant" of the universal vacuum.
\end{itemize}


\section{Evolution across Cosmological Epochs}
The mechanical interpretation of $G$ implies a relationship with the condensate density $\rho_m$ that may vary across cosmological time. In a tracking regime where $\rho_m$ follows the background density evolution, the gravitational coupling $G$ scales inversely with the field density.

\begin{table}[h]
\centering
\caption{Evolution of $G$ and $\rho_m$ (Hypothetical Tracking Regime)}
\label{tab:epochs}
\vspace{2mm}
\begin{tabular}{@{}llll@{}}
\toprule
\textbf{Epoch} & \textbf{Redshift ($z$)} & \textbf{Rel. Density $\rho_m / \rho_0$} & \textbf{Rel. Coupling $G/G_0$} \\ \midrule
Radiation Dom. & $> 3000$ & $\sim 10^{10}$ & $\sim 10^{-10}$ \\
Matter Dom. & $\sim 1$ & $8$ & $0.125$ \\
Present Day & $0$ & $1$ & $1$ \\ \bottomrule
\end{tabular}
\end{table}

\section{Phenomenological Implications}
The variable nature of $G$ in this framework offers a potential resolution to several contemporary cosmological tensions. Notably, the high-redshift massive galaxies recently observed by the \textit{James Webb Space Telescope} (JWST) suggest a faster growth of structure than predicted by standard $\Lambda$CDM.\cite{BoylanKolchin2023} A stronger gravitational coupling $G$ in the early universe—driven by a lower initial condensate density $\rho_m$ in certain models—may provide the necessary acceleration for such early assembly. 


Furthermore, the mechanical interpretation of $G$ impacts black hole physics; in a denser, "stiffer" vacuum (higher $\rho_m$, lower $G$), Hawking radiation temperatures $T_H \propto (GM)^{-1}$ would be suppressed, altering evaporation rates. Finally, this Condensate Link Framework may address challenges in Fuzzy Dark Matter (FDM) by linking gravitational strength to the local field pressure, potentially resolving the "cusp-core" problem by softening the potential where the condensate density is highest.



\section{Conclusion}
In this work, we have demonstrated that Newton's gravitational constant $G$ can be derived from the mechanical bulk modulus of a dark-matter vacuum condensate. By identifying the vacuum as a barotropic fluid with energy density $u_\nu$, we matched the hydrodynamic field pressure to the geometric stiffness of the Einstein--Hilbert action. This mechanical route not only recovers the observed value of $G$ through cosmological observables but may provide a dynamic framework that accounts for anomalies in early structure formation observed by JWST and structural tensions in dark matter halos. Our results suggest that gravity is an emergent phenomenon, with $G$ functioning as the physical stiffness of the universal medium.

\section*{Declarations}
\subsection*{AI Transparency Statement}
The fundamental theoretical framework, the novel correlation of matching mechanical bulk modulus to geometric stiffness, and the conceptual derivation of $G$ are the original work of the author. Generative AI (Google Gemini 3 Flash) was utilized specifically for LaTeX document formatting, bibliography organization, research and assistance with mathematical formulas, notation and scientific prose to meet journal standards.

\section*{Acknowledgments}
\subsection{Novel Correlation and Synthesis of Established Frameworks}
The derivation presented in this work is a synthesis of several established pillars of modern physics. We acknowledge and express profound gratitude to the foundational contributions of the scientific community including, but not limited to; the quantum condensate community, particularly the work on the Gross--Pitaevskii equation, the pioneers of analogue gravity, and the teams responsible for the Planck and JWST missions. Their pursuit of precision and their willingness to explore the boundaries of physics provided the data and theories required to anchor this mechanical route. Furthermore, this derivation relies on the heat kernel methods developed within the quantum field theory community and the high-precision measurements of cosmological parameters provided by the global astrophysical community.


The unique contribution of this research is the \textbf{correlation} of these disparate frameworks through a specific mechanical-to-geometric matching condition. By demonstrating that the bulk modulus of the dark-matter medium is identically equal to the Einstein--Hilbert prefactor, we provide a new physical derivation for the value of $G$. This transforms $G$ from an empirical constant into a calculable property, honoring the insights of these various fields while providing a unified mechanical explanation for the strength of gravity.


This research did not receive any specific grant from funding agencies in the public, commercial, or not-for-profit sectors, and remains an independent contribution to the field.

\appendix
\section{One--loop derivation of the EH term}
\label{app:one-loop}

We integrate out the field fluctuations $\phi$. The divergent part of the one-loop effective action $\Gamma_{1}$ is given by the Schwinger-proper time representation:\cite{Vassilevich03,Schwinger1951}
\begin{equation}
\Gamma_{1} = -\frac{1}{2(4\pi)^2} \int d^4x \sqrt{-g} \int_{1/\Lambda_{UV}^2}^\infty \frac{ds}{s^3} e^{-sM_{\rm eff}^2} [a_0 + a_1 s + a_2 s^2].
\end{equation}
The term $a_1 = (\tfrac{1}{6}-\xi)R$ induces the EH action:
\begin{equation}
S_{\rm ind} \supset \frac{(\tfrac{1}{6}-\xi)\Lambda_{UV}^2}{32\pi^2} \int d^4x \sqrt{-g} R.
\end{equation}
Matching this to $S_{\rm EH} = \frac{c^4}{16\pi G} \int \sqrt{-g} R$ anchors $G$ to the medium's cutoff $\Lambda_{UV}$.

\begin{thebibliography}{99}
\bibitem{Unruh1981} W. G. Unruh, Phys. Rev. Lett. \textbf{46}, 1351 (1981).
\bibitem{Visser1998} M. Visser, Class. Quantum Grav. \textbf{15}, 1767 (1998).
\bibitem{Barcelo2011} C. Barcel\'{o}, S. Liberati, and M. Visser, Living Rev. Relativ. \textbf{14}, 3 (2011).
\bibitem{Einstein1916} A. Einstein, Annalen der Physik \textbf{354}, 769 (1916).
\bibitem{Hilbert1915} D. Hilbert, Konigl. Gesell. d. Wiss. Göttingen, Nachr. Math.-Phys. Kl. 395 (1915).
\bibitem{Gross1961} E. P. Gross, Nuovo Cimento \textbf{20}, 454 (1961).
\bibitem{Pitaevskii1961} L. P. Pitaevskii, Sov. Phys. JETP \textbf{13}, 451 (1961).
\bibitem{LandauLifshitz} L. D. Landau and E. M. Lifshitz, \textit{Fluid Mechanics}, 2nd ed. (Pergamon, London, 1987).
\bibitem{Vassilevich03} D. V. Vassilevich, Phys. Rep. \textbf{388}, 279 (2003).
\bibitem{Schwinger1951} J. Schwinger, Phys. Rev. \textbf{82}, 664 (1951).
\bibitem{PDG2023} O. Lahav and A. R. Liddle, in \textit{Review of Particle Physics}, PTEP \textbf{2022}, 083C01 (2022).
\bibitem{BoylanKolchin2023} M. Boylan-Kolchin, Nature Astron. \textbf{7}, 731 (2023).
\end{thebibliography}

\end{document}